
\section{Application: Exact community recovery}
\label{sec:community-detection-linf}


Another application that benefits remarkably from the fine-grained eigenvector perturbation theory is community recovery. This section reexplores the stochastic block model studied in Section~\ref{sec:community-detection}, and develops significantly enhanced theoretical support for spectral clustering. 




\subsection{Performance guarantees: Exact recovery}
\label{sec:performance-exact-recovery}


The focus of this section is simultaneous recovery of the community memberships of all vertices, 
which is termed {\em exact recovery} or {\em strong consistency} in the community detection literature \citep{abbe2017community}. This imposes a much stronger requirement than the {\em weak consistency} studied in Section~\ref{sec:theory-community-detection}. 

For the sake of conciseness, the theorem below concentrates on the challenging regime where $p,q\asymp \frac{\log n}{n}$,  corresponding to the lowest possible edge densities that allow for exact recovery. This is because, if $p < \frac{\log n}{n}$, then with high probability, one can find isolated vertices that are not connected with any edge in the graph \citep{durrett2007random}; hence, there will be absolutely no means to infer the community membership of these isolated vertices.  The theoretical guarantee is as follows. 
%
\begin{theorem}
	\label{thm:community-recovery-linf}
	Fix any constant $\varepsilon>0$, and consider the setting of Section~\ref{sec:setup-community-detection}. Suppose  $p=\frac{\alpha \log n}{n}$ and $q=\frac{\beta \log n}{n}$ for some sufficiently large constants $\alpha > \beta >0$.\footnote{In the current proof, the constants $\alpha$, $\beta$ might depend on the fixed choice of $\varepsilon$. Encouragingly, this restriction can be lifted; see~\citet{abbe2020entrywise} for details.} In addition, assume that
	%
	\begin{align}
		\big( \sqrt{p} -\sqrt{q}\big)^2 \geq   2\left( 1+ \varepsilon  \right) \frac{\log n}{n} .
		\label{eq:H-pq-lower-bound-theorem}
	\end{align}
	%
	With probability  $1-o(1)$, the spectral method in Section~\ref{sec:algorithm-community-detection} yields
	%
	\[
		x_{i}=x_{i}^{\star}~~\text{for all }1\leq i\leq n, \quad~\text{or}\quad~ x_{i}=-x_{i}^{\star}~~\text{for all }1\leq i\leq n.
	\]
\end{theorem}
%
This theorem, which first appeared in \citet{abbe2020entrywise},  identifies a sufficient recovery condition in terms of the edge densities. 
The result substantially strengthens the $\ell_2$-based theory in Section~\ref{sec:theory-community-detection}, 
uncovering the capability of the spectral method in achieving not merely almost exact recovery in the average sense, 
but more appealingly,  exact community recovery that ensures correct labels of all vertices.  

A natural question arises as to whether the recovery condition \eqref{eq:H-pq-lower-bound-theorem} is improvable via more sophisticated algorithms. Answering this question requires information-theoretic thinking, that is, how to characterize a fundamental threshold---in terms of the difference of edge densities---below which exact recovery is deemed  infeasible. As has been demonstrated in \citet{abbe2014exact,mossel2015consistency,hajek2015achieving}, no algorithm whatsoever is able to achieve exact community recovery  if
%
\begin{equation}
	\big( \sqrt{p} -\sqrt{q}\big)^2 \leq  2(1-\varepsilon) \frac{\log n}{n} 
	\label{eq:IT-lower-bound-CD}
\end{equation}
%
for any constant $\varepsilon >0$. 
This fundamental lower bound, in conjunction with Theorem~\ref{thm:community-recovery-linf}, reveals a sharp phase transition behind the performance of the spectral method. In particular, its optimality is guaranteed all the way down to the information-theoretic threshold; see Figure~\ref{fig:sbm-inf} for numerical evidence.  

\begin{figure}[!t]
\begin{center}
	\includegraphics[width=0.55\textwidth]{figures/SBM-rec-new.pdf} 
\end{center}	
	\caption{Phase transition of spectral methods for exact community recovery. Set $n=300$, $p=(a \log n) / n$, and $q = (b \log n) / n$. We vary $a,b$ from 0 to 8 with an equal space of 0.1. For each configuration of $a \geq b$, we conduct 100 Monte Carlo trials and report the empirical success rate for recovering the entire community structure correctly. 
	%White areas indicate complete success, while dark ones indicate complete failure. 
	The empirical phase transition occurs near the information-theoretic threshold (see \eqref{eq:IT-lower-bound-CD}).
	\label{fig:sbm-inf}}
\end{figure}



Given that the above information-theoretic threshold is specified in terms of $( \sqrt{p} -\sqrt{q})^2$,
the reader might naturally wonder what the operational meaning of this quantity is. As it turns out, this metric is a sort of distance measure between the two edge probability distributions under consideration. In truth, in the setting of Theorem~\ref{thm:community-recovery-linf}, this metric is intimately related to the squared Hellinger distance between two Bernoulli distributions.
%
\begin{definition}[Squared Hellinger distance]
%
\label{defn:Hellinger-distance}
Consider two distributions $P$ and $Q$ over a finite alphabet $\mathcal{Y}$. The squared Hellinger distance $\mathsf{H}^2(P\,\|\,Q)$ between $P$ and $Q$  is defined as follows	%
\begin{equation}
\mathsf{H}^2(P\,\|\,Q)\coloneqq\frac{1}{2}\sum\nolimits_{y\in \mathcal{Y}}\Big(\sqrt{P(y)}-\sqrt{Q(y)}\Big)^{2}.
\label{eq:defn-Hellinger-PQ}
\end{equation}
%
\end{definition}
%
In particular, consider the squared Hellinger distance between two Bernoulli distributions of interest $\mathsf{Bern}(p)$ and $\mathsf{Bern}(q)$, where we denote by $\mathsf{Bern}(p)$ the Bernoulli distribution with mean $p$. It is seen that \citep{chen2016information}
%
\begin{align*}
	\mathsf{H}^2\big( \mathsf{Bern}(p),   \mathsf{Bern}(q) \big) 
	&\coloneqq \frac{1}{2} \big( \sqrt{p} -\sqrt{q}\big)^2 + \frac{1}{2} \big( \sqrt{1-p} -\sqrt{1-q}\big)^2 \nonumber\\
	&= (1+o(1)) \frac{1}{2} \big( \sqrt{p} -\sqrt{q}\big)^2,
\end{align*}
%
when $p=o(1)$ and $q=o(1)$.\footnote{To justify this approximation, the following calculation suffices: \[ \sqrt{1-q}-\sqrt{1-p}=\frac{p-q}{\sqrt{1-p}+\sqrt{1-q}}=(1+o(1))\big(\sqrt{p}-\sqrt{q}\big)\big(\sqrt{p}+\sqrt{q}\big)=o\big(\sqrt{p}-\sqrt{q}\big).\]} 
The phase transition phenomenon identified in \eqref{eq:H-pq-lower-bound-theorem} and \eqref{eq:IT-lower-bound-CD} can then be alternatively described as
%
\begin{align*}
\text{\text{spectral method works}}\quad & \text{if }\mathsf{H}^{2}\big(\mathsf{Bern}(p),\mathsf{Bern}(q)\big)\geq(1+\varepsilon)\frac{\log n}{n}\\
\text{no algorithm works}\quad & \text{if }\mathsf{H}^{2}\big(\mathsf{Bern}(p),\mathsf{Bern}(q)\big)\leq(1-\varepsilon)\frac{\log n}{n}
\end{align*}
%
for an arbitrary small constant $\varepsilon >0$. 


\subsection{Proof of Theorem~\ref{thm:community-recovery-linf}} 

We now turn to the proof of Theorem~\ref{thm:community-recovery-linf}. 
Without loss of generality, suppose that $x_i^{\star}=1$ for all $1\leq i\leq n/2$  and $x_i^{\star}=-1$ for all $i> n/2$, so that 
$\bm{u}^{\star} =\frac{1}{\sqrt{n}}  {\small
\left[\begin{array}{c}
\bm{1}_{n/2}\\
-\bm{1}_{n/2}
\end{array}\right] }$. 



Recalling the matrix $\bM$ given in \eqref{eq:M-data-matrix-SBM} and its mean $\bM^{\star}$ in \eqref{eq:M-data-matrix-expectation-SBM}, 
one can immediately see that $\kappa=\mu=r=1$ for $\bM^{\star}$ in this application. 
Theorem \ref{thm:UsgnH-Ustar-MUstar-general} (cf.~\eqref{eq:UsgnH-MUstar-bound-theorem-general}) readily implies 
the existence of some $z\in\{1,-1\}$ such that
%
\begin{align}
 & \Big\| z\bm{u}-\frac{1}{\lambda^{\star}}\bm{M}\bm{u}^{\star}\Big\|_{\infty}
\lesssim\frac{\sigma}{|\lambda^{\star}|}+\frac{\sigma^{2}\sqrt{n\log n}+\sigma B \,{\log^{3/2}n}}{(\lambda^{\star})^{2}} 
\label{eq:implication-Thm-symmetry-1st-CD}
\end{align}
%
with probability at least $1-O(n^{-5})$. Additionally, it has already been explained in Section~\ref{sec:community-detection} that  
%
\[
	 B=1,\quad\sigma^{2}\leq\max\{p,q\}=p, \quad \text{and} \quad \lambda^{\star}={n(p-q)}/{2}.
\]
%
Substitution into \eqref{eq:implication-Thm-symmetry-1st-CD} reveals that
%
\begin{align}
	\big\| z\lambda^{\star}\bm{u}-\bm{M}\bm{u}^{\star} \big\|_{\infty} & \lesssim\sigma+\frac{\sigma^{2}\sqrt{n\log n}}{\lambda^{\star}}+\frac{\sigma B\,\log^{3/2}n}{\lambda^{\star}}\nonumber \\
 	& \leq  C \Big( \sqrt{p}+\frac{p\sqrt{\log n}}{\sqrt{n}(p-q)}+\frac{\sqrt{p}\log^{3/2}n}{n(p-q)} \Big)
	\label{eq:z-lambda-u-MU-bound-CD}
\end{align}
%
holds for some universal constant $C>0$.
As a result, a crucial step boils down to controlling $\bm{M}\bm{u}^{\star}$ in an entrywise manner: each element is a difference between two independent random binomial random variables and  is accomplished through the following lemma. 
%
\begin{lemma}
	\label{lemma:M-ustar-lower-bound-CD}
	Suppose that
%
\begin{align}
	\mathsf{H}_{p,q}^{2} \coloneqq \big( \sqrt{p} -\sqrt{q}\big)^2 \geq\left( 1+ \varepsilon  \right) \frac{2\log n}{n} 
	\label{eq:H-pq-lower-bound-lemma}
\end{align}
% 
for some quantity $\varepsilon>0$.  Let $\varepsilon_0 \coloneqq \frac{\varepsilon\log n}{\sqrt{n}\log\frac{p(1-q)}{q(1-p)}}-\frac{1}{\sqrt{n}}$. 
Then with probability exceeding $1-n^{-\varepsilon/2}$, one has
%
\begin{align*}
	 \bm{M}_{l,\cdot}\bm{u}^{\star}   \geq \varepsilon_0 \,\,\,\text{for all }l\leq \frac{n}{2},
	\quad\text{and}\quad
	   \bm{M}_{l,\cdot} \bm{u}^{\star}  
	\leq - \varepsilon_0 \,\,\, \text{for all } l > \frac{n}{2}. 
	%\label{eq:M-ustar-bound-CD-lemma}
\end{align*}
%
\end{lemma}
%

We now return to analyze the entrywise behavior of $\bm{u}$. Note that
%
\begin{align*}
(z\lambda^{\star})u_{l} & \geq\bm{M}_{l,\cdot}\bm{u}^{\star}-\big| z\lambda^{\star}u_{l}-\bm{M}_{l,\cdot}\bm{u}^{\star}\big|  \quad\text{for all } l\leq n/2;\\
(z\lambda^{\star})u_{l} & \leq\bm{M}_{l,\cdot}\bm{u}^{\star}+\big| z\lambda^{\star}u_{l}-\bm{M}_{l,\cdot}\bm{u}^{\star}\big|  \quad\text{for all }l>n/2.
\end{align*}
%
This together with \eqref{eq:z-lambda-u-MU-bound-CD}, Lemma~\ref{lemma:M-ustar-lower-bound-CD} and $\lambda^{\star}>0$ yields that if
%
\begin{align}
	\frac{\varepsilon\log n}{\sqrt{n}\log\frac{p(1-q)}{q(1-p)}}  > \frac{1}{\sqrt{n}} 
	+  C \Big( \sqrt{p} + \frac{p\sqrt{\log n}}{\sqrt{n}(p-q)} + \frac{\sqrt{p}\log^{3/2}n}{n(p-q)} \Big), 
	\label{eq:epsilon-condition-CD-inf-123}
\end{align}
%
then it follows that
%
\[
zu_{l}>0\quad\text{for all }1\leq l\leq\frac{n}{2},\quad\text{and}\quad zu_{l}<0\quad\text{for all }l>\frac{n}{2},
\]
%
thus guaranteeing exact community recovery once the rounding procedure (based on the sign) is applied.  


To finish up, it remains to validate Condition~\eqref{eq:epsilon-condition-CD-inf-123}. 
Fixing $\varepsilon>0$ to be a constant, we make the following observations. 
%
\begin{itemize}
	\item From the assumptions $\varepsilon \asymp 1$ and $p,q\asymp \frac{\log n}{n}$ (or $\alpha,\beta \asymp 1$), one has
	%
	\[
		\frac{\varepsilon\log n}{\sqrt{n}\log\frac{p(1-q)}{q(1-p)}}=\frac{\varepsilon\log n}{\sqrt{n}\log\frac{(1+o(1))\alpha}{\beta}}
		\asymp \frac{\log n}{\sqrt{n}}\gg \sqrt{p}+\frac{1}{\sqrt{n}}.
	\]
	\item Turning to the term $\frac{p\sqrt{\log n}}{\sqrt{n}(p-q)}$,   we  observe that 
	%
	\begin{align}
		& \log\frac{p(1-q)}{q(1-p)}=\log\Big(1+\frac{p-q}{q(1-p)}\Big)\leq\frac{p-q}{q(1-p)}\leq\frac{2(\alpha-\beta)}{\beta}, \notag\\
		& \qquad\quad \Longrightarrow \quad 
		\frac{\varepsilon\log n}{\sqrt{n}\log\frac{p(1-q)}{q(1-p)}}\geq\frac{\varepsilon\log n}{2\sqrt{n}}\,\frac{\beta}{\alpha-\beta},
		\label{eq:LB-CD-12345}
	\end{align}
	%
	where in the last inequality of the  first line we have used the assumption $p=o(1)$ and hence $1/(1-p) \leq 2$. 
	Given that $\varepsilon, \alpha,\beta \asymp 1$, it is guaranteed that
	%
	\[
		\eqref{eq:LB-CD-12345} \gg \frac{\alpha\sqrt{\log n}}{\sqrt{n}(\alpha-\beta)}=\frac{p\sqrt{\log n}}{\sqrt{n}(p-q)} .
	\]
	% 
	\item We then move on to the term $\frac{\sqrt{p}\log^{3/2}n}{n(p-q)}$. 
	If $\alpha / \beta \leq 2$ and  $\beta \geq\frac{200C^{2}}{\varepsilon^{2}} \geq \frac{100C^{2}\alpha/\beta}{\varepsilon^{2}}$, then one  has $\beta \geq \frac{10C\sqrt{\alpha}}{\varepsilon}$ and hence
	%
	\[
		\eqref{eq:LB-CD-12345} \geq\frac{5C\sqrt{\alpha}\log n}{\sqrt{n}(\alpha-\beta)}=\frac{5C\sqrt{p}\log^{3/2}n}{n(p-q)} .
	\]
	%
	 In addition, if $\alpha / \beta > 2$, then it follows that
	%
	\begin{equation}
		\frac{\sqrt{p}\log^{3/2}n}{n(p-q)}=\frac{\sqrt{\alpha}\log n}{\sqrt{n}(\alpha-\beta)}\leq\frac{2\sqrt{\alpha}\log n}{\sqrt{n}\alpha}=\frac{2\log n}{\sqrt{n\alpha}},
		\label{eq:LB-CD-5678}
	\end{equation}
	%
	where the inequality holds true since $\alpha - \beta > \alpha - \alpha/2 = \alpha / 2$. 
	Using the basic inequality $\log x \leq \sqrt{x}$ further leads  to
	%
	\[
		\frac{\varepsilon\log n}{\sqrt{n}\log\frac{p(1-q)}{q(1-p)}}\geq\frac{\varepsilon\log n}{\sqrt{n}\log\frac{2\alpha}{\beta}}
		\geq\frac{\varepsilon\sqrt{\beta}\log n}{\sqrt{n}\sqrt{2\alpha}}\geq\frac{5C\sqrt{p}\log^{3/2}n}{n(p-q)}.
	\]
	%
	Here, the first inequality holds since $p,q=o(1)$ and hence $\frac{1-q}{1-p}\leq 2$, 
	whereas the last relation relies on \eqref{eq:LB-CD-5678} and holds with the proviso that $\beta\geq200C^{2}/\varepsilon^{2}$. 
	
 
\end{itemize}
%
The above calculations taken collectively establish Condition~\eqref{eq:epsilon-condition-CD-inf-123} under the assumptions of Theorem~\ref{thm:community-recovery-linf}, thus concluding the proof. 



\begin{remark}
	It is worth pointing out that the bound \eqref{eq:UsgnH-Ustar-bound-theorem-general} in Theorem \ref{thm:UsgnH-Ustar-MUstar-general} is not sufficiently tight when establishing this result. 
	Instead, one needs to resort to the more refined bound \eqref{eq:UsgnH-MUstar-bound-theorem-general} in Theorem \ref{thm:UsgnH-Ustar-MUstar-general}, 
	which allows us to sharpen the error bound by explicitly accounting for the first-order error term $(\bm{M}-\bm{M}^{\star})\bm{u}^{\star}$. 
\end{remark}




 


\subsection{Proof of auxiliary lemmas}
\label{sec:proof-auxiliary-CD}

Before embarking on the proof of Lemma~\ref{lemma:M-ustar-lower-bound-CD}, we first record non-asymptotic tail bounds concerning log-likelihood ratios and a sum of Bernoulli random variables, which make apparent the role of the squared Hellinger distance \citep{tsybakov2009nonparm}.
%
\begin{lemma}
	\label{lemma:LLR-Hellinger}
	Consider two distributions $P$ and $Q$ over a finite alphabet $\mathcal{Y}$, and suppose that $P(y)\neq 0$ for all $y\in \mathcal{Y}$. Generate an independent sequence $\{y_i\}_{1\leq i\leq n}$ obeying $y_i \sim P$. Then for any $\zeta\in \mathbb{R}$ one has
	%
\begin{equation}
\mathbb{P}\left\{ \sum_{i=1}^{n}\log\frac{Q(y_{i})}{P(y_{i})}\geq-n\zeta\right\} \leq\exp\left(-n\Big[\mathsf{H}^2(P\,\|\,Q)-\frac{\zeta}{2}\Big]\right),
\label{eq:prob-log-PQ-exp-Hel}
\end{equation}
%
where $\mathsf{H}^2(P\,\|\,Q)$ is the squared Hellinger distance between $P$ and $Q$ defined in \eqref{eq:defn-Hellinger-PQ}.
\end{lemma}
%
%
\begin{lemma}
\label{lem:sum-Bernoulli-CD}
%
Consider two sequences of independent random variables
%
\[
z_{i}\sim\mathsf{Bern}(p),\qquad w_{i}\sim\mathsf{Bern}(q),\qquad1\leq i\leq n,
\]
%
and suppose that $p>q$. For any $\xi\in \mathbb{R}$, it follows that
%
\begin{align*}
\mathbb{P}\left\{ \sum_{i=1}^{n}(z_{i}-w_i) \leq n\xi\right\}  & \leq\exp\left(-n\Big[\mathsf{H}_{p,q}^{2}-\frac{\xi}{2}\log\frac{p(1-q)}{q(1-p)}\Big]\right) ,
\end{align*}
%
where $\mathsf{H}_{p,q}^{2}  \coloneqq  \big(\sqrt{p}-\sqrt{q} \, \big)^{2}$.
%
%\begin{align}
%	.
%	\label{eq:defn-Hpq}
%\end{align}
% 
\end{lemma}
%
In what follows, we first establish Lemmas~\ref{lemma:LLR-Hellinger} and \ref{lem:sum-Bernoulli-CD}, and then return to prove Lemma~\ref{lemma:M-ustar-lower-bound-CD}. 



\paragraph{Proof of Lemma~\ref{lemma:LLR-Hellinger}.}
Apply the Chernoff bound to yield 
%
\begin{align}
 & \mathbb{P}\left\{ \sum_{i=1}^{n}\log\frac{Q(y_{i})}{P(y_{i})}\geq-n\zeta\right\} {\leq}\frac{\mathbb{E}_{y_{i}\sim P}\left[\exp\left(\frac{1}{2}\sum_{i=1}^{n}\log\frac{Q(y_{i})}{P(y_{i})}\right)\right]}{\exp(-n\zeta/2)} \notag\\
 & \qquad \qquad\overset{(\mathrm{i})}{=}
 %\frac{\prod_{i=1}^{n}\mathbb{E}_{y_{i}\sim P}\left[\exp\left(\frac{1}{2}\log\frac{Q(y_{i})}{P(y_{i})}\right)\right]}{\exp(-n\zeta/2)}
 %= 
 \frac{\left(\mathbb{E}_{y\sim P}\Big[\left(\frac{Q(y)}{P(y)}\right)^{1/2}\Big]\right)^{n}}{\exp(-n\zeta/2)},
\label{eq:prob-log-PQ-UB}
\end{align}
%
where 
%(i) results from the Chernoff bound, and (ii) 
(i) holds due to the i.i.d.~assumption of the $y_{i}$'s.  In addition,
%It is seen from the definition of $P$ and $Q$ that
%
\begin{align}
\mathbb{E}_{y\sim P}\left[\left(\frac{Q(y)}{P(y)}\right)^{1/2}\right] & =\sum_{y}P(y)\left(\frac{Q(y)}{P(y)}\right)^{1/2}=\sum_{y}\sqrt{P(y)Q(y)} \notag\\
 & =1-\frac{1}{2}\sum_{y}\left(P(y)+Q(y)-2\sqrt{P(y)Q(y)}\right) \notag\\
  & =1-\frac{1}{2}\sum_{y}\left( \sqrt{P(y)} - \sqrt{Q(y)} \right)^2 \notag\\
 & =1-\mathsf{H}^2(P\,\|\,Q)\leq  \exp \big(-\mathsf{H}^2(P\,\|\,Q) \big),
 \label{eq:P-Q-Hellinger-bound}
\end{align}
%
where the second line follows since $\sum_yP(y) = \sum_yQ(y)=1$, and  the last line uses the definition \eqref{eq:defn-Hellinger-PQ} and the elementary inequality $1-x\leq \exp(-x)$. 
Substituting \eqref{eq:P-Q-Hellinger-bound} into \eqref{eq:prob-log-PQ-UB} concludes the proof. 







\paragraph{Proof of Lemma~\ref{lem:sum-Bernoulli-CD}.}

Set  $y_i\coloneqq z_i - w_i$. The proof is built upon a mapping between $\sum_{i=1}^{n}y_{i}$ and a certain log-likelihood ratio. 
Specifically, let us introduce two distributions $P$ and $Q$ supported on $\{1,0,-1\}$:
%
\[
P(x)=\begin{cases}
p(1-q),\quad & \text{if }x=1,\\
q(1-p), & \text{if }x=-1,\\
pq+(1-p)(1-q), \quad & \text{if }x=0,
\end{cases} 
%\qquad \text{and}
\]
% 
%
\[
Q(x)=\begin{cases}
q(1-p),\quad & \text{if }x=1,\\
p(1-q), & \text{if }x=-1,\\
pq+(1-p)(1-q), \quad & \text{if }x=0.
\end{cases}
\]
%
Apparently, $P$ (resp.~$Q$) corresponds to the distribution of $y_i$ (resp.~$-y_i$).  A key observation is that
%
\begin{align*}
\sum_{i=1}^{n}\log\frac{Q(y_{i})}{P(y_{i})} & =\sum_{i=1}^{n}\left\{ \mathbbm1\{y_{i}=1\}\log\frac{q(1-p)}{p(1-q)}+\mathbbm1\{y_{i}=-1\}\log\frac{p(1-q)}{q(1-p)}\right\} \\
 & =\sum_{i=1}^{n}y_{i}\log\frac{q(1-p)}{p(1-q)} ,
\end{align*}
%
which relies on the fact that $y_i$ is supported on $\{1,0,-1\}$. 
Recognizing that $\log\frac{p(1-q)}{q(1-p)}>0$ holds as long as
$p>q$ (since $q(1-p)<p(1-q)$), we can further derive
%
\begin{align*}
\mathbb{P}\left\{ \sum_{i=1}^{n}y_{i}\leq n\xi\right\}  & =\mathbb{P}\left\{ \frac{1}{\log\frac{p(1-q)}{q(1-p)}}\sum_{i=1}^{n}\log\frac{Q(y_{i})}{P(y_{i})}\geq-n\xi\right\} \\
 & \leq\exp\left(-n\Big[\mathsf{H}^2(P\,\|\,Q)-\frac{\xi}{2}\log\frac{p(1-q)}{q(1-p)}\Big]\right),
\end{align*}
%
where the last inequality comes from Lemma~\ref{lemma:LLR-Hellinger}. 
From the constructions of $P$ and $Q$ and the definition \eqref{eq:defn-Hellinger-PQ} of $\mathsf{H}^2(P\,\|\,Q)$, it is easily seen that 
%
\begin{align*}
\mathsf{H}^{2}(P\,\|\,Q) & =\Big(\sqrt{p(1-q)}-\sqrt{q(1-p)}\,\Big)^{2}
	= \frac{\big(p-q\big)^{2}}{\big(\sqrt{p(1-q)}+\sqrt{q(1-p)}\,\big)^{2}}\\
 & \geq\frac{\big(p-q\big)^{2}}{\big(\sqrt{p}+\sqrt{q}\,\big)^{2}}=\big(\sqrt{p}-\sqrt{q}\big)^{2} ,
\end{align*}
%
thus concluding the proof. 






\paragraph{Proof of Lemma~\ref{lemma:M-ustar-lower-bound-CD}.}

Let us start by looking at the first entry of $\bm{M}\bm{u}^{\star}$. 
It is seen from the construction \eqref{eq:M-data-matrix-SBM} that
%
\begin{align}
\bm{M}_{1,\cdot}\bm{u}^{\star}=\bm{A}_{1,\cdot}\bm{u}^{\star}-\frac{p+q}{2}\big(\bm{1}^{\top}\bm{u}^{\star}\big)\bm{1}+pu_{1}^{\star}
\geq \bm{A}_{1,\cdot}\bm{u}^{\star} ,
%+ O\Big(\frac{1}{\sqrt{n}}\Big),
	\label{eq:connection-M1u-A1u-CD}
\end{align}
%
where we have used the fact that $\bm{1}^{\top}\bm{u}^{\star}=0$  and $u_1^{\star}>0$. 
The expression  $\bm{u}^{\star} =\frac{1}{\sqrt{n}}  {\small
\left[\begin{array}{c}
\bm{1}_{n/2}\\
-\bm{1}_{n/2}
\end{array}\right] }$ admits the following decomposition
%
\begin{align}
	\bm{A}_{1,\cdot}\bm{u}^{\star}
%=\frac{1}{\sqrt{n}}\sum_{i=2}^{n/2}\big(A_{1,i}-A_{1,i+n/2}\big)+\frac{1}{\sqrt{n}}A_{1,n/2+1}
	= \frac{1}{\sqrt{n}}\sum\nolimits_{i=1}^{n/2}\big(A_{1,i}-A_{1,i+n/2}\big) , 
	\label{eq:A1ustar-decompose-CD}
\end{align}
%
which can be controlled via Lemma~\ref{lem:sum-Bernoulli-CD}. 
%Here and throughout, we denote by $\mathsf{Bern}(p)$  the Bernoulli distribution with mean $p$.



Observe that $A_{1,i} \sim \mathsf{Bern}(p)$ for all $1< i\leq n/2$ and $A_{1,i} \sim \mathsf{Bern}(q)$ otherwise. Using the definitions of $z_i$ and $w_i$ in Lemma~\ref{lem:sum-Bernoulli-CD}, we obtain
%
\begin{align}
 & \mathbb{P}\left\{ \sum\nolimits_{i=1}^{n/2}\big(A_{1,i}-A_{1,i+n/2}\big)\leq\frac{n\zeta}{2}-1\right\}  
	\leq\mathbb{P}\left\{ \sum\nolimits_{i=1}^{n/2}\big(z_{i}-w_{i}\big)\leq\frac{n\zeta}{2}\right\}  \notag\\
	& \qquad \leq\exp\left(-\frac{n}{2}\Big[\mathsf{H}_{p,q}^{2}-\frac{\zeta}{2}\log\frac{p(1-q)}{q(1-p)}\Big]\right) \leq \frac{1}{n^{1+\delta}}
	\label{eq:A-sum-complicated-CD}
\end{align}
%
for some $\delta >0$, where the first inequality follows since $A_{1,1}=0\leq z_1+1$ (so that $A_{1,1}-A_{1,n/2+1}$ is stochastically dominated by $z_1 - w_1$), and the last inequality holds as long as
%
\begin{equation}
	\mathsf{H}_{p,q}^{2}-\frac{\zeta}{2}\log\frac{p(1-q)}{q(1-p)}
	\geq\frac{2(1+\delta)\log n}{n} ,
	\label{eq:H-pq-lower-bound-123}
\end{equation}
which we shall ensure at the end of the proof. Substituting \eqref{eq:A-sum-complicated-CD} into \eqref{eq:A1ustar-decompose-CD} and \eqref{eq:connection-M1u-A1u-CD} yields
\begin{align*}
\mathbb{P}\left\{ \bm{M}_{1,\cdot}\bm{u}^{\star}\leq \frac{n\zeta-2}{2\sqrt{n}} \right\}  
%& \leq\mathbb{P}\left\{ \bm{A}_{1,\cdot}\bm{u}^{\star}\leq\varepsilon_{0}-1/\sqrt{n}\right\} \\
 & \leq \mathbb{P}\left\{ \bm{A}_{1,\cdot}\bm{u}^{\star}\leq\frac{n\zeta-2}{2\sqrt{n}} \right\} \leq\frac{1}{n^{1+\delta}}.
\end{align*}
%
%which combined with \eqref{eq:connection-M1u-A1u-CD} leads to
% and the assumption ${1}/{\lambda^{\star}} = o(\varepsilon_0)$ leads to
%
%\[
%	\mathbb{P}\left\{ \bm{M}_{1,\cdot}\bm{u}^{\star} \leq  {\varepsilon_0}  \right\} \leq
%	\mathbb{P}\left\{ \bm{A}_{1,\cdot}\bm{u}^{\star} \leq  {\varepsilon_0}   \right\} 
%	\leq\frac{1}{n^{1+\delta}}.
%\]
%
Repeating the preceding analysis for $ \bm{M}_{l,\cdot}\bm{u}^{\star}$ with other  $l$'s  and taking the union bound, we see that with probability at least $1- n^{-\delta}$,
%
\begin{subequations}
	\label{eq:M-ustar-bound-CD}
\begin{align}
%\begin{split}
% & \frac{1}{\lambda^{\star}}\bm{A}_{l,\cdot}\bm{u}^{\star}\geq\frac{\sqrt{n}\zeta}{4\lambda^{\star}}\ \qquad\,\text{for all }l\leq\frac{n}{2}\\
% & \frac{1}{\lambda^{\star}}\bm{A}_{l,\cdot}\bm{u}^{\star}\leq-\frac{\sqrt{n}\zeta}{4\lambda^{\star}}\ \ \quad \text{for all }l>\frac{n}{2}
%\end{split}
	\bm{M}_{l,\cdot}\bm{u}^{\star} &\geq \frac{n\zeta-2}{2\sqrt{n}}
	%\frac{\sqrt{n}\zeta}{4\lambda^{\star}}
	\quad  &&\text{if } l\leq {n}/{2} \\
	%\quad
	%\text{and}\quad
	\bm{M}_{l,\cdot}\bm{u}^{\star}
	&\leq - \frac{n\zeta-2}{2\sqrt{n}}
	%-\frac{\sqrt{n}\zeta}{4\lambda^{\star}}
	\quad  &&\text{if } l> n/2
	%\big)
\end{align}
\end{subequations}
%
%\begin{equation}
%\begin{array}{cc}
%	& \frac{1}{\lambda^{\star}}\bm{M}_{l,\cdot}\bm{u}^{\star}  \geq\frac{\sqrt{n}\zeta-2}{2\lambda^{\star}},\ \quad~~\text{}l\leq\frac{n}{2}\\
%	& \frac{1}{\lambda^{\star}}\bm{M}_{l,\cdot}\bm{u}^{\star}  \leq-\frac{\sqrt{n}\zeta-2}{2\lambda^{\star}},\ \ ~~\text{}l>\frac{n}{2}
%\end{array}
%\end{equation}
%
hold simultaneously for all $1\leq l\leq n$. 



Finally, it remains to ensure satisfaction of \eqref{eq:H-pq-lower-bound-123}. As it turns out, if the condition \eqref{eq:H-pq-lower-bound-lemma}
holds, then it suffices to take $\delta\leq \varepsilon / 2$ and 
$\zeta=\frac{2\varepsilon\log n}{n\log\frac{p(1-q)}{q(1-p)}}$. This completes the proof.  


